% pracovni_list.tex -- Lekce 08: Funkce
\newcommand{\lekce}{Lekce 08 -- Pracovní list}
% preamble-pl.tex -- sdílená preamble pro pracovní listy
% Includuje se z ucitel/lekceXX/pracovni_list.tex jako:
%   % preamble-pl.tex -- sdílená preamble pro pracovní listy
% Includuje se z ucitel/lekceXX/pracovni_list.tex jako:
%   % preamble-pl.tex -- sdílená preamble pro pracovní listy
% Includuje se z ucitel/lekceXX/pracovni_list.tex jako:
%   \input{../spolecne/preamble-pl.tex}

\documentclass[a4paper,12pt]{article}

% --- Jazyk a font ---
\usepackage{polyglossia}
\setmainlanguage{czech}
\usepackage{fontspec}

% --- Okraje ---
\usepackage[top=20mm, bottom=20mm, left=22mm, right=22mm]{geometry}

% --- Česky korektní uvozovky ---
\usepackage{csquotes}
\newcommand{\uv}[1]{\enquote{#1}}

% --- Typografie ---
\usepackage{microtype}
\usepackage{parskip}          % odstavce odděleny mezerou, ne odsazením

% --- Barvy a grafika ---
\usepackage[table]{xcolor}
\definecolor{microbit-blue}{HTML}{1E4D8C}
\definecolor{code-bg}{HTML}{F5F5F5}
\definecolor{task-bg}{HTML}{EEF4FF}

% --- Tabulky ---
\usepackage{booktabs}
\usepackage{multirow}

% --- Zdrojový kód (Python) ---
\usepackage{listings}
\lstset{
  language=Python,
  basicstyle=\ttfamily\small,
  backgroundcolor=\color{code-bg},
  frame=single,
  framerule=0pt,
  rulecolor=\color{gray!30},
  xleftmargin=6pt,
  xrightmargin=6pt,
  breaklines=true,
  showstringspaces=false,
  keywordstyle=\color{microbit-blue}\bfseries,
  stringstyle=\color{teal},
  commentstyle=\color{gray}\itshape,
  literate=
    {á}{{\'a}}1 {é}{{\'e}}1 {í}{{\'i}}1 {ó}{{\'o}}1 {ú}{{\'u}}1
    {ů}{{\r{u}}}1 {č}{{\v{c}}}1 {š}{{\v{s}}}1 {ž}{{\v{z}}}1
    {Á}{{\'A}}1 {É}{{\'E}}1 {Í}{{\'I}}1 {Ó}{{\'O}}1 {Ú}{{\'U}}1
    {Č}{{\v{C}}}1 {Š}{{\v{S}}}1 {Ž}{{\v{Z}}}1 {ň}{{\v{n}}}1 {ř}{{\v{r}}}1,
}

% --- Zadání úkolů ---
\usepackage[most]{tcolorbox}
\tcbuselibrary{breakable}

\newcounter{ukol}[section]
\newenvironment{ukol}[1][]{%
  \stepcounter{ukol}%
  \begin{tcolorbox}[
    colback=task-bg,
    colframe=microbit-blue,
    fonttitle=\bfseries,
    title={Úkol~\theukol\ifx\relax#1\relax\else{: #1}\fi},
    breakable,
    left=6pt, right=6pt, top=4pt, bottom=4pt
  ]%
}{%
  \end{tcolorbox}%
}

% Inline kód — underscore package zajistí, ze _ funguje v texttt
\usepackage{underscore}
\newcommand{\pyinline}[1]{\texttt{#1}}
% Pojem (zvýraznění termínu) — stejné jako v preamble-prez.tex
\newcommand{\pojem}[1]{\textbf{#1}}

% Místo pro odpověď studenta
\newcommand{\odpoved}[1][3cm]{%
  \vspace{#1}%
}

% --- Záhlaví / zápatí ---
\usepackage{fancyhdr}
\pagestyle{fancy}
\fancyhf{}
\renewcommand{\headrulewidth}{0.4pt}
\fancyhead[L]{\small\textcolor{microbit-blue}{\textbf{Úvod do Pythonu na Micro:bitu}}}
\fancyhead[R]{\small\textcolor{gray}{\lekce}}
\fancyfoot[C]{\small\thepage}

% --- Ostatní ---
\usepackage{enumitem}
\usepackage{graphicx}
\usepackage{hyperref}
\hypersetup{
  colorlinks=true,
  linkcolor=microbit-blue,
  urlcolor=microbit-blue,
}


\documentclass[a4paper,12pt]{article}

% --- Jazyk a font ---
\usepackage{polyglossia}
\setmainlanguage{czech}
\usepackage{fontspec}

% --- Okraje ---
\usepackage[top=20mm, bottom=20mm, left=22mm, right=22mm]{geometry}

% --- Česky korektní uvozovky ---
\usepackage{csquotes}
\newcommand{\uv}[1]{\enquote{#1}}

% --- Typografie ---
\usepackage{microtype}
\usepackage{parskip}          % odstavce odděleny mezerou, ne odsazením

% --- Barvy a grafika ---
\usepackage[table]{xcolor}
\definecolor{microbit-blue}{HTML}{1E4D8C}
\definecolor{code-bg}{HTML}{F5F5F5}
\definecolor{task-bg}{HTML}{EEF4FF}

% --- Tabulky ---
\usepackage{booktabs}
\usepackage{multirow}

% --- Zdrojový kód (Python) ---
\usepackage{listings}
\lstset{
  language=Python,
  basicstyle=\ttfamily\small,
  backgroundcolor=\color{code-bg},
  frame=single,
  framerule=0pt,
  rulecolor=\color{gray!30},
  xleftmargin=6pt,
  xrightmargin=6pt,
  breaklines=true,
  showstringspaces=false,
  keywordstyle=\color{microbit-blue}\bfseries,
  stringstyle=\color{teal},
  commentstyle=\color{gray}\itshape,
  literate=
    {á}{{\'a}}1 {é}{{\'e}}1 {í}{{\'i}}1 {ó}{{\'o}}1 {ú}{{\'u}}1
    {ů}{{\r{u}}}1 {č}{{\v{c}}}1 {š}{{\v{s}}}1 {ž}{{\v{z}}}1
    {Á}{{\'A}}1 {É}{{\'E}}1 {Í}{{\'I}}1 {Ó}{{\'O}}1 {Ú}{{\'U}}1
    {Č}{{\v{C}}}1 {Š}{{\v{S}}}1 {Ž}{{\v{Z}}}1 {ň}{{\v{n}}}1 {ř}{{\v{r}}}1,
}

% --- Zadání úkolů ---
\usepackage[most]{tcolorbox}
\tcbuselibrary{breakable}

\newcounter{ukol}[section]
\newenvironment{ukol}[1][]{%
  \stepcounter{ukol}%
  \begin{tcolorbox}[
    colback=task-bg,
    colframe=microbit-blue,
    fonttitle=\bfseries,
    title={Úkol~\theukol\ifx\relax#1\relax\else{: #1}\fi},
    breakable,
    left=6pt, right=6pt, top=4pt, bottom=4pt
  ]%
}{%
  \end{tcolorbox}%
}

% Inline kód — underscore package zajistí, ze _ funguje v texttt
\usepackage{underscore}
\newcommand{\pyinline}[1]{\texttt{#1}}
% Pojem (zvýraznění termínu) — stejné jako v preamble-prez.tex
\newcommand{\pojem}[1]{\textbf{#1}}

% Místo pro odpověď studenta
\newcommand{\odpoved}[1][3cm]{%
  \vspace{#1}%
}

% --- Záhlaví / zápatí ---
\usepackage{fancyhdr}
\pagestyle{fancy}
\fancyhf{}
\renewcommand{\headrulewidth}{0.4pt}
\fancyhead[L]{\small\textcolor{microbit-blue}{\textbf{Úvod do Pythonu na Micro:bitu}}}
\fancyhead[R]{\small\textcolor{gray}{\lekce}}
\fancyfoot[C]{\small\thepage}

% --- Ostatní ---
\usepackage{enumitem}
\usepackage{graphicx}
\usepackage{hyperref}
\hypersetup{
  colorlinks=true,
  linkcolor=microbit-blue,
  urlcolor=microbit-blue,
}


\documentclass[a4paper,12pt]{article}

% --- Jazyk a font ---
\usepackage{polyglossia}
\setmainlanguage{czech}
\usepackage{fontspec}

% --- Okraje ---
\usepackage[top=20mm, bottom=20mm, left=22mm, right=22mm]{geometry}

% --- Česky korektní uvozovky ---
\usepackage{csquotes}
\newcommand{\uv}[1]{\enquote{#1}}

% --- Typografie ---
\usepackage{microtype}
\usepackage{parskip}          % odstavce odděleny mezerou, ne odsazením

% --- Barvy a grafika ---
\usepackage[table]{xcolor}
\definecolor{microbit-blue}{HTML}{1E4D8C}
\definecolor{code-bg}{HTML}{F5F5F5}
\definecolor{task-bg}{HTML}{EEF4FF}

% --- Tabulky ---
\usepackage{booktabs}
\usepackage{multirow}

% --- Zdrojový kód (Python) ---
\usepackage{listings}
\lstset{
  language=Python,
  basicstyle=\ttfamily\small,
  backgroundcolor=\color{code-bg},
  frame=single,
  framerule=0pt,
  rulecolor=\color{gray!30},
  xleftmargin=6pt,
  xrightmargin=6pt,
  breaklines=true,
  showstringspaces=false,
  keywordstyle=\color{microbit-blue}\bfseries,
  stringstyle=\color{teal},
  commentstyle=\color{gray}\itshape,
  literate=
    {á}{{\'a}}1 {é}{{\'e}}1 {í}{{\'i}}1 {ó}{{\'o}}1 {ú}{{\'u}}1
    {ů}{{\r{u}}}1 {č}{{\v{c}}}1 {š}{{\v{s}}}1 {ž}{{\v{z}}}1
    {Á}{{\'A}}1 {É}{{\'E}}1 {Í}{{\'I}}1 {Ó}{{\'O}}1 {Ú}{{\'U}}1
    {Č}{{\v{C}}}1 {Š}{{\v{S}}}1 {Ž}{{\v{Z}}}1 {ň}{{\v{n}}}1 {ř}{{\v{r}}}1,
}

% --- Zadání úkolů ---
\usepackage[most]{tcolorbox}
\tcbuselibrary{breakable}

\newcounter{ukol}[section]
\newenvironment{ukol}[1][]{%
  \stepcounter{ukol}%
  \begin{tcolorbox}[
    colback=task-bg,
    colframe=microbit-blue,
    fonttitle=\bfseries,
    title={Úkol~\theukol\ifx\relax#1\relax\else{: #1}\fi},
    breakable,
    left=6pt, right=6pt, top=4pt, bottom=4pt
  ]%
}{%
  \end{tcolorbox}%
}

% Inline kód — underscore package zajistí, ze _ funguje v texttt
\usepackage{underscore}
\newcommand{\pyinline}[1]{\texttt{#1}}
% Pojem (zvýraznění termínu) — stejné jako v preamble-prez.tex
\newcommand{\pojem}[1]{\textbf{#1}}

% Místo pro odpověď studenta
\newcommand{\odpoved}[1][3cm]{%
  \vspace{#1}%
}

% --- Záhlaví / zápatí ---
\usepackage{fancyhdr}
\pagestyle{fancy}
\fancyhf{}
\renewcommand{\headrulewidth}{0.4pt}
\fancyhead[L]{\small\textcolor{microbit-blue}{\textbf{Úvod do Pythonu na Micro:bitu}}}
\fancyhead[R]{\small\textcolor{gray}{\lekce}}
\fancyfoot[C]{\small\thepage}

% --- Ostatní ---
\usepackage{enumitem}
\usepackage{graphicx}
\usepackage{hyperref}
\hypersetup{
  colorlinks=true,
  linkcolor=microbit-blue,
  urlcolor=microbit-blue,
}


\begin{document}

\begin{center}
  {\Large\bfseries\textcolor{microbit-blue}{Lekce 08: Funkce}}\\[4pt]
  {\large Pracovní list}
\end{center}

\medskip\noindent
\pojem{Funkce} zabalí opakující se kód do pojmenovaného celku.
Definujeme ji jednou pomocí \pyinline{def} a voláme libovolněkrát.

% ------------------------------------------------------------------
\begin{ukol}[Funkce pozdrav]

Napíšeme funkci, která zobrazí pozdrav se jménem.

\begin{lstlisting}
from microbit import *

def pozdrav(jmeno):
    display.scroll("Ahoj, " + jmeno + "!")

pozdrav("Karel")
pozdrav("Jana")
pozdrav("Svete")
\end{lstlisting}

\begin{enumerate}
  \item Zadej program a spusť ho — co se zobrazí?
  \item Co je \pyinline{jmeno} — parametr nebo argument? A co \pyinline{"Karel"}?
  \item Uprav funkci tak, aby zobrazila i \pyinline{Image.HAPPY} po pozdravení.
  \item Zkus zavolat \pyinline{pozdrav} bez argumentu: \pyinline{pozdrav()}.
        Jakou chybu Python ohlásí?
\end{enumerate}

\odpoved[2cm]

\begin{tcolorbox}[colback=code-bg, colframe=gray!40,
                  fonttitle=\small\bfseries, title=Klíčové pojmy]
  \small
  \begin{itemize}[leftmargin=*]
    \item \pyinline{def nazev(parametr):} — definuje funkci (neprovede ji!).
    \item \pyinline{nazev(argument)} — zavolá funkci, spustí tělo.
    \item \pyinline{return hodnota} — ukončí funkci a vrátí výsledek.
    \item Funkce bez \pyinline{return} vrátí \pyinline{None}.
  \end{itemize}
\end{tcolorbox}

\end{ukol}

% ------------------------------------------------------------------
\begin{ukol}[Hod kostkou s \texttt{return}]

Přepíšeme hod kostkou z lekce~05 jako funkci s návratovou hodnotou.

\begin{lstlisting}
from microbit import *
import random

def hod_kostkou(pocet_sten):
    return random.randint(1, pocet_sten)

while True:
    if accelerometer.was_gesture("shake"):
        cislo = hod_kostkou(6)
        display.show(str(cislo))
    sleep(100)
\end{lstlisting}

\begin{enumerate}
  \item Zadej program a otestuj ho zatřesením.
  \item Co by se stalo, kdybys zapomněl \pyinline{return}?
        Zkus to — co se zobrazí na displeji?
  \item Uprav program tak, aby tlačítko A hodilo šestistěnnou kostku
        a tlačítko B dvacetistěnnou (1--20).
\end{enumerate}

\odpoved[2cm]

\end{ukol}

% ------------------------------------------------------------------
\begin{ukol}[Refaktoring]

Vezmi libovolný program z~předchozích lekcí (teploměr, věštírna, jukebox\ldots)
a přepiš ho pomocí funkcí.

\medskip
Podmínky:
\begin{itemize}
  \item definuj alespoň \textbf{2 vlastní funkce}
  \item alespoň jedna funkce musí mít \textbf{parametr}
  \item alespoň jedna funkce musí \textbf{vracet hodnotu} (\pyinline{return})
  \item hlavní smyčka \pyinline{while True:} by měla být \textbf{přehledná a krátká}
\end{itemize}

\medskip
\textit{Tip: Začni tím, že identifikuješ opakující se bloky kódu — ty se stanou funkcemi.}

\odpoved[6cm]

\noindent Výsledný program ulož a připrav k odevzdání do Google Classroom.

\end{ukol}

% ------------------------------------------------------------------
\begin{bonusukol}[Funkce vracející seznam]

Funkce může vracet nejen číslo, ale i celý seznam výsledků:

\begin{lstlisting}
from microbit import *
import random

def generuj_hody(pocet, strany=6):
    vysledky = []
    for i in range(pocet):
        vysledky.append(random.randint(1, strany))
    return vysledky

while True:
    if accelerometer.was_gesture("shake"):
        hody  = generuj_hody(3)        # tri sestistenne kostky
        soucet = 0
        for h in hody:
            soucet = soucet + h
        display.scroll(str(soucet))
    sleep(100)
\end{lstlisting}

\begin{enumerate}
  \item Zadej program a zatřes — zobrazí se součet tří kostek.
  \item Jaký typ dat vrací \pyinline{generuj_hody()}? Jak to poznáš?
  \item Uprav volání tak, aby házel 5 dvacetistěnných kostek:
        \pyinline{generuj_hody(5, strany=20)}.
  \item \textbf{Výzva:} Napiš funkci \pyinline{prumer(seznam)}, která vrátí
        průměr čísel ze seznamu (\pyinline{sum(seznam) // len(seznam)}).
        Zobraz průměr hodu spolu se součtem.
\end{enumerate}

\odpoved[3cm]

\end{bonusukol}

\end{document}
